% stub.tex -- course-neutral stand-in for a real exam template, so that question
% fragments can be compiled and checked on their own. Same macro names as the real
% templates (templates/examples/*.sty); plain layout. \input it after \documentclass.
% Needs XeLaTeX or LuaLaTeX (fontspec), e.g. tectonic -X compile standalone.tex.
%
%   \documentclass[11pt,a4paper]{article}
%   % stub.tex -- course-neutral stand-in for a real exam template, so that question
% fragments can be compiled and checked on their own. Same macro names as the real
% templates (templates/examples/*.sty); plain layout. \input it after \documentclass.
% Needs XeLaTeX or LuaLaTeX (fontspec), e.g. tectonic -X compile standalone.tex.
%
%   \documentclass[11pt,a4paper]{article}
%   % stub.tex -- course-neutral stand-in for a real exam template, so that question
% fragments can be compiled and checked on their own. Same macro names as the real
% templates (templates/examples/*.sty); plain layout. \input it after \documentclass.
% Needs XeLaTeX or LuaLaTeX (fontspec), e.g. tectonic -X compile standalone.tex.
%
%   \documentclass[11pt,a4paper]{article}
%   % stub.tex -- course-neutral stand-in for a real exam template, so that question
% fragments can be compiled and checked on their own. Same macro names as the real
% templates (templates/examples/*.sty); plain layout. \input it after \documentclass.
% Needs XeLaTeX or LuaLaTeX (fontspec), e.g. tectonic -X compile standalone.tex.
%
%   \documentclass[11pt,a4paper]{article}
%   \input{stub.tex}
%   \begin{document} \input{question.tex} \end{document}
\usepackage[a4paper,margin=2.5cm]{geometry}
\usepackage{fontspec}
% \setmainfont{Arial}   % optional: the font of your real paper, if it is installed
\usepackage{amsmath,amssymb}
\usepackage[american,siunitx]{circuitikz}   % loads tikz and siunitx
\usetikzlibrary{arrows.meta,calc}
\usepackage{enumitem}
\setlength{\parindent}{0pt}
\setlength{\parskip}{6pt}
\pagestyle{plain}

% notation: bold vectors, sans-serif components and matrices, sans-serif transpose
\providecommand{\vect}[1]{\mathbf{#1}}
\providecommand{\mat}[1]{\mathsf{#1}}
\providecommand{\T}{^{\mathsf{T}}}

\makeatletter
% \begin{Question}{<number>}{<total marks>} ... \end{Question}: new page, heading, closing line.
% \begin{Question}[continue]{..}{..} sets it on the same page (layout tests only).
\newenvironment{Question}[3][]{%
  \def\stub@tmp{#1}\def\stub@cont{continue}%
  \ifx\stub@tmp\stub@cont\par\bigskip\else\clearpage\fi
  \gdef\stub@qnum{#2}\gdef\stub@qtotal{#3}%
  \noindent\textbf{Question #2}\par}
 {\par\nobreak\noindent\hfill\textbf{[End of Question \stub@qnum; \stub@qtotal\ marks total]}\par}
% parts (a), (b): \begin{Parts} \item ... \end{Parts};  sub-parts (i), (ii): \begin{SubParts} \item ... \end{SubParts}
\newlist{Parts}{enumerate}{1}
\setlist[Parts]{label=(\alph*), leftmargin=2.4em, itemsep=6pt}
\newlist{SubParts}{enumerate}{1}
\setlist[SubParts]{label=(\roman*), leftmargin=2.4em, itemsep=4pt}
% \Marks{9}, \Marks{4, 5}, \Mark{1}: "[9 marks]" flush right at the end of the line
% (on a line of its own when used between paragraphs); \Marks{1} prints "[1 mark]"
\def\stub@one{1}
\def\stub@putmarks#1{%
  \ifvmode\noindent\hfill[#1]\par
  \else{\unskip\nobreak\hfil\penalty50\hskip2em\hbox{}\nobreak\hfil[#1]%
    \parfillskip=0pt\finalhyphendemerits=0\par}\fi}
\newcommand{\Marks}[1]{\edef\stub@tmp{\zap@space#1 \@empty}%
  \ifx\stub@tmp\stub@one\stub@putmarks{#1 mark}\else\stub@putmarks{#1 marks}\fi}
\newcommand{\Mark}[1]{\stub@putmarks{#1 mark}}
% \ExamFigure{<tikz/circuitikz picture>}{Q1.1} or \ExamFigure[<caption text>]{..}{Q1.1}
\newcommand{\ExamFigure}[3][]{\par\begin{center}#2\par\vspace{4pt}%
  Figure #3\if\relax\detokenize{#1}\relax\else: #1\fi\end{center}\par}
% \AnswerSpace{<marks>}: blank, unbreakable writing space, 1.6 cm per mark, at least 4 cm,
% at most 95% of the text height; "4, 5" is summed
\newlength{\stub@ans}
\newcommand{\AnswerSpace}[1]{\par\setlength{\stub@ans}{0pt}%
  \@for\stub@m:=#1\do{\addtolength{\stub@ans}{\stub@m\dimexpr 1.6cm\relax}}%
  \ifdim\stub@ans<4cm\setlength{\stub@ans}{4cm}\fi
  \ifdim\stub@ans>0.95\textheight\setlength{\stub@ans}{0.95\textheight}\fi
  \noindent\rule{0pt}{\stub@ans}\par}
% \EndOfPaper{<text>}: closing line after the last question
\newcommand{\EndOfPaper}[1]{\par\vspace{2em}{\centering\textbf{#1}\par}}
\makeatother

%   \begin{document} \input{question.tex} \end{document}
\usepackage[a4paper,margin=2.5cm]{geometry}
\usepackage{fontspec}
% \setmainfont{Arial}   % optional: the font of your real paper, if it is installed
\usepackage{amsmath,amssymb}
\usepackage[american,siunitx]{circuitikz}   % loads tikz and siunitx
\usetikzlibrary{arrows.meta,calc}
\usepackage{enumitem}
\setlength{\parindent}{0pt}
\setlength{\parskip}{6pt}
\pagestyle{plain}

% notation: bold vectors, sans-serif components and matrices, sans-serif transpose
\providecommand{\vect}[1]{\mathbf{#1}}
\providecommand{\mat}[1]{\mathsf{#1}}
\providecommand{\T}{^{\mathsf{T}}}

\makeatletter
% \begin{Question}{<number>}{<total marks>} ... \end{Question}: new page, heading, closing line.
% \begin{Question}[continue]{..}{..} sets it on the same page (layout tests only).
\newenvironment{Question}[3][]{%
  \def\stub@tmp{#1}\def\stub@cont{continue}%
  \ifx\stub@tmp\stub@cont\par\bigskip\else\clearpage\fi
  \gdef\stub@qnum{#2}\gdef\stub@qtotal{#3}%
  \noindent\textbf{Question #2}\par}
 {\par\nobreak\noindent\hfill\textbf{[End of Question \stub@qnum; \stub@qtotal\ marks total]}\par}
% parts (a), (b): \begin{Parts} \item ... \end{Parts};  sub-parts (i), (ii): \begin{SubParts} \item ... \end{SubParts}
\newlist{Parts}{enumerate}{1}
\setlist[Parts]{label=(\alph*), leftmargin=2.4em, itemsep=6pt}
\newlist{SubParts}{enumerate}{1}
\setlist[SubParts]{label=(\roman*), leftmargin=2.4em, itemsep=4pt}
% \Marks{9}, \Marks{4, 5}, \Mark{1}: "[9 marks]" flush right at the end of the line
% (on a line of its own when used between paragraphs); \Marks{1} prints "[1 mark]"
\def\stub@one{1}
\def\stub@putmarks#1{%
  \ifvmode\noindent\hfill[#1]\par
  \else{\unskip\nobreak\hfil\penalty50\hskip2em\hbox{}\nobreak\hfil[#1]%
    \parfillskip=0pt\finalhyphendemerits=0\par}\fi}
\newcommand{\Marks}[1]{\edef\stub@tmp{\zap@space#1 \@empty}%
  \ifx\stub@tmp\stub@one\stub@putmarks{#1 mark}\else\stub@putmarks{#1 marks}\fi}
\newcommand{\Mark}[1]{\stub@putmarks{#1 mark}}
% \ExamFigure{<tikz/circuitikz picture>}{Q1.1} or \ExamFigure[<caption text>]{..}{Q1.1}
\newcommand{\ExamFigure}[3][]{\par\begin{center}#2\par\vspace{4pt}%
  Figure #3\if\relax\detokenize{#1}\relax\else: #1\fi\end{center}\par}
% \AnswerSpace{<marks>}: blank, unbreakable writing space, 1.6 cm per mark, at least 4 cm,
% at most 95% of the text height; "4, 5" is summed
\newlength{\stub@ans}
\newcommand{\AnswerSpace}[1]{\par\setlength{\stub@ans}{0pt}%
  \@for\stub@m:=#1\do{\addtolength{\stub@ans}{\stub@m\dimexpr 1.6cm\relax}}%
  \ifdim\stub@ans<4cm\setlength{\stub@ans}{4cm}\fi
  \ifdim\stub@ans>0.95\textheight\setlength{\stub@ans}{0.95\textheight}\fi
  \noindent\rule{0pt}{\stub@ans}\par}
% \EndOfPaper{<text>}: closing line after the last question
\newcommand{\EndOfPaper}[1]{\par\vspace{2em}{\centering\textbf{#1}\par}}
\makeatother

%   \begin{document} \input{question.tex} \end{document}
\usepackage[a4paper,margin=2.5cm]{geometry}
\usepackage{fontspec}
% \setmainfont{Arial}   % optional: the font of your real paper, if it is installed
\usepackage{amsmath,amssymb}
\usepackage[american,siunitx]{circuitikz}   % loads tikz and siunitx
\usetikzlibrary{arrows.meta,calc}
\usepackage{enumitem}
\setlength{\parindent}{0pt}
\setlength{\parskip}{6pt}
\pagestyle{plain}

% notation: bold vectors, sans-serif components and matrices, sans-serif transpose
\providecommand{\vect}[1]{\mathbf{#1}}
\providecommand{\mat}[1]{\mathsf{#1}}
\providecommand{\T}{^{\mathsf{T}}}

\makeatletter
% \begin{Question}{<number>}{<total marks>} ... \end{Question}: new page, heading, closing line.
% \begin{Question}[continue]{..}{..} sets it on the same page (layout tests only).
\newenvironment{Question}[3][]{%
  \def\stub@tmp{#1}\def\stub@cont{continue}%
  \ifx\stub@tmp\stub@cont\par\bigskip\else\clearpage\fi
  \gdef\stub@qnum{#2}\gdef\stub@qtotal{#3}%
  \noindent\textbf{Question #2}\par}
 {\par\nobreak\noindent\hfill\textbf{[End of Question \stub@qnum; \stub@qtotal\ marks total]}\par}
% parts (a), (b): \begin{Parts} \item ... \end{Parts};  sub-parts (i), (ii): \begin{SubParts} \item ... \end{SubParts}
\newlist{Parts}{enumerate}{1}
\setlist[Parts]{label=(\alph*), leftmargin=2.4em, itemsep=6pt}
\newlist{SubParts}{enumerate}{1}
\setlist[SubParts]{label=(\roman*), leftmargin=2.4em, itemsep=4pt}
% \Marks{9}, \Marks{4, 5}, \Mark{1}: "[9 marks]" flush right at the end of the line
% (on a line of its own when used between paragraphs); \Marks{1} prints "[1 mark]"
\def\stub@one{1}
\def\stub@putmarks#1{%
  \ifvmode\noindent\hfill[#1]\par
  \else{\unskip\nobreak\hfil\penalty50\hskip2em\hbox{}\nobreak\hfil[#1]%
    \parfillskip=0pt\finalhyphendemerits=0\par}\fi}
\newcommand{\Marks}[1]{\edef\stub@tmp{\zap@space#1 \@empty}%
  \ifx\stub@tmp\stub@one\stub@putmarks{#1 mark}\else\stub@putmarks{#1 marks}\fi}
\newcommand{\Mark}[1]{\stub@putmarks{#1 mark}}
% \ExamFigure{<tikz/circuitikz picture>}{Q1.1} or \ExamFigure[<caption text>]{..}{Q1.1}
\newcommand{\ExamFigure}[3][]{\par\begin{center}#2\par\vspace{4pt}%
  Figure #3\if\relax\detokenize{#1}\relax\else: #1\fi\end{center}\par}
% \AnswerSpace{<marks>}: blank, unbreakable writing space, 1.6 cm per mark, at least 4 cm,
% at most 95% of the text height; "4, 5" is summed
\newlength{\stub@ans}
\newcommand{\AnswerSpace}[1]{\par\setlength{\stub@ans}{0pt}%
  \@for\stub@m:=#1\do{\addtolength{\stub@ans}{\stub@m\dimexpr 1.6cm\relax}}%
  \ifdim\stub@ans<4cm\setlength{\stub@ans}{4cm}\fi
  \ifdim\stub@ans>0.95\textheight\setlength{\stub@ans}{0.95\textheight}\fi
  \noindent\rule{0pt}{\stub@ans}\par}
% \EndOfPaper{<text>}: closing line after the last question
\newcommand{\EndOfPaper}[1]{\par\vspace{2em}{\centering\textbf{#1}\par}}
\makeatother

%   \begin{document} \input{question.tex} \end{document}
\usepackage[a4paper,margin=2.5cm]{geometry}
\usepackage{fontspec}
% \setmainfont{Arial}   % optional: the font of your real paper, if it is installed
\usepackage{amsmath,amssymb}
\usepackage[american,siunitx]{circuitikz}   % loads tikz and siunitx
\usetikzlibrary{arrows.meta,calc}
\usepackage{enumitem}
\setlength{\parindent}{0pt}
\setlength{\parskip}{6pt}
\pagestyle{plain}

% notation: bold vectors, sans-serif components and matrices, sans-serif transpose
\providecommand{\vect}[1]{\mathbf{#1}}
\providecommand{\mat}[1]{\mathsf{#1}}
\providecommand{\T}{^{\mathsf{T}}}

\makeatletter
% \begin{Question}{<number>}{<total marks>} ... \end{Question}: new page, heading, closing line.
% \begin{Question}[continue]{..}{..} sets it on the same page (layout tests only).
\newenvironment{Question}[3][]{%
  \def\stub@tmp{#1}\def\stub@cont{continue}%
  \ifx\stub@tmp\stub@cont\par\bigskip\else\clearpage\fi
  \gdef\stub@qnum{#2}\gdef\stub@qtotal{#3}%
  \noindent\textbf{Question #2}\par}
 {\par\nobreak\noindent\hfill\textbf{[End of Question \stub@qnum; \stub@qtotal\ marks total]}\par}
% parts (a), (b): \begin{Parts} \item ... \end{Parts};  sub-parts (i), (ii): \begin{SubParts} \item ... \end{SubParts}
\newlist{Parts}{enumerate}{1}
\setlist[Parts]{label=(\alph*), leftmargin=2.4em, itemsep=6pt}
\newlist{SubParts}{enumerate}{1}
\setlist[SubParts]{label=(\roman*), leftmargin=2.4em, itemsep=4pt}
% \Marks{9}, \Marks{4, 5}, \Mark{1}: "[9 marks]" flush right at the end of the line
% (on a line of its own when used between paragraphs); \Marks{1} prints "[1 mark]"
\def\stub@one{1}
\def\stub@putmarks#1{%
  \ifvmode\noindent\hfill[#1]\par
  \else{\unskip\nobreak\hfil\penalty50\hskip2em\hbox{}\nobreak\hfil[#1]%
    \parfillskip=0pt\finalhyphendemerits=0\par}\fi}
\newcommand{\Marks}[1]{\edef\stub@tmp{\zap@space#1 \@empty}%
  \ifx\stub@tmp\stub@one\stub@putmarks{#1 mark}\else\stub@putmarks{#1 marks}\fi}
\newcommand{\Mark}[1]{\stub@putmarks{#1 mark}}
% \ExamFigure{<tikz/circuitikz picture>}{Q1.1} or \ExamFigure[<caption text>]{..}{Q1.1}
\newcommand{\ExamFigure}[3][]{\par\begin{center}#2\par\vspace{4pt}%
  Figure #3\if\relax\detokenize{#1}\relax\else: #1\fi\end{center}\par}
% \AnswerSpace{<marks>}: blank, unbreakable writing space, 1.6 cm per mark, at least 4 cm,
% at most 95% of the text height; "4, 5" is summed
\newlength{\stub@ans}
\newcommand{\AnswerSpace}[1]{\par\setlength{\stub@ans}{0pt}%
  \@for\stub@m:=#1\do{\addtolength{\stub@ans}{\stub@m\dimexpr 1.6cm\relax}}%
  \ifdim\stub@ans<4cm\setlength{\stub@ans}{4cm}\fi
  \ifdim\stub@ans>0.95\textheight\setlength{\stub@ans}{0.95\textheight}\fi
  \noindent\rule{0pt}{\stub@ans}\par}
% \EndOfPaper{<text>}: closing line after the last question
\newcommand{\EndOfPaper}[1]{\par\vspace{2em}{\centering\textbf{#1}\par}}
\makeatother
